% ============================================================
\chapter{Epidemiological Models}
\label{ch:epidemiology}
% ============================================================

\section{Introduction}

In 1927, William Ogilvy Kermack and Anderson Gray McKendrick published a paper that would found mathematical epidemiology: their SIR model divides the population into three compartments---Susceptible, Infectious, Recovered---and predicts, with remarkable simplicity, the course of an epidemic. The basic reproduction number $R_0$ emerges naturally: if $R_0 > 1$, the epidemic spreads; if $R_0 < 1$, it dies out. This threshold, which became famous during the COVID-19 pandemic, is a purely mathematical result that guides public health policy.

Compartmental epidemiological models allow us to understand the spread
of infectious diseases and evaluate the impact of interventions
(vaccination, quarantine).  This chapter develops the SIR and SEIR
models from first principles, analyses the basic reproduction number
$\mathcal{R}_0$, and provides Python simulations.

% ------------------------------------------------------------
\section{The Kermack--McKendrick SIR Model}
\label{sec:sir}
% ------------------------------------------------------------

\begin{definition}[SIR compartments]
The total population $N$ (constant) is divided into three compartments:
\begin{itemize}
  \item $S(t)$: \textbf{Susceptible} (individuals who can be infected).
  \item $I(t)$: \textbf{Infectious} (contagious individuals).
  \item $R(t)$: \textbf{Recovered} (immunised or deceased).
\end{itemize}
Conservation constraint: $S(t)+I(t)+R(t)=N$.
\end{definition}

\subsection{Model Derivation}

\begin{intuition}[title=\textbf{SIR assumptions}]
\begin{itemize}
  \item The population is homogeneous and well-mixed.
  \item The transmission rate is proportional to the product $SI$
    (mass action).
  \item Infectious individuals recover at constant rate $\gamma$.
  \item No births, natural deaths, or demography.
\end{itemize}
\end{intuition}

The SIR equations are:
\begin{align}
  \frac{dS}{dt} &= -\beta\frac{SI}{N}, \label{eq:sir_s}\\
  \frac{dI}{dt} &= \beta\frac{SI}{N} - \gamma I, \label{eq:sir_i}\\
  \frac{dR}{dt} &= \gamma I, \label{eq:sir_r}
\end{align}
where $\beta>0$ is the transmission rate and $\gamma>0$ the recovery
rate.

\begin{keyformulas}[title=\textbf{SIR model parameters}]
\begin{center}
\begin{tabular}{ccc}
\toprule
Parameter & Meaning & Unit \\
\midrule
$\beta$ & Transmission rate & $\text{time}^{-1}$ \\
$\gamma$ & Recovery rate & $\text{time}^{-1}$ \\
$1/\gamma$ & Mean infectious period & time \\
$\mathcal{R}_0 = \beta/\gamma$ & Basic reproduction number & dimensionless \\
\bottomrule
\end{tabular}
\end{center}
\end{keyformulas}

% ------------------------------------------------------------
\section{The Basic Reproduction Number $\mathcal{R}_0$}
\label{sec:r0}
% ------------------------------------------------------------

\begin{definition}[Basic reproduction number]
The \textbf{basic reproduction number} $\mathcal{R}_0$ is the average
number of secondary cases produced by one infectious individual in a
fully susceptible population:
\[
  \mathcal{R}_0 = \frac{\beta}{\gamma}.
\]
\end{definition}

\begin{theorem}[Epidemic threshold]
An epidemic occurs if and only if $\mathcal{R}_0 > 1$.  More
precisely, $dI/dt > 0$ at the start of the epidemic if and only if
$\mathcal{R}_0 S(0)/N > 1$.
\end{theorem}

\begin{proof}
At the start, $S\approx N$.  Equation~\eqref{eq:sir_i} gives:
\[
  \frac{dI}{dt}\bigg|_{t=0} \approx (\beta-\gamma)I_0
  = \gamma(\mathcal{R}_0-1)I_0.
\]
Thus $dI/dt > 0$ if and only if $\mathcal{R}_0 > 1$.
\end{proof}

\begin{remark}
Typical values of $\mathcal{R}_0$: measles $\approx 12$--$18$,
influenza $\approx 1.5$--$3$, COVID-19 (original strain)
$\approx 2.5$--$3.5$, Ebola $\approx 1.5$--$2.5$.
\end{remark}

% ------------------------------------------------------------
\section{Qualitative Analysis of the SIR Model}
% ------------------------------------------------------------

\subsection{Reduction to Two Equations}

Since $R=N-S-I$, it suffices to study the $(S,I)$ system in the
triangle $\{S\geq 0,\;I\geq 0,\;S+I\leq N\}$.

\subsection{Final Size of the Epidemic}

\begin{theorem}[Final size equation]
If $I(0)=I_0\ll N$ and $S(0)\approx N$, the final fraction of
susceptibles $s_\infty = S(\infty)/N$ satisfies the transcendental
equation:
\[
  \ln s_\infty = \mathcal{R}_0(s_\infty - 1).
\]
\end{theorem}

\begin{proof}
Dividing $dS/dI$:
\[
  \frac{dS}{dI} = \frac{-\beta SI/N}{\beta SI/N - \gamma I}
  = \frac{-\beta S/N}{\beta S/N - \gamma}.
\]
Setting $s = S/N$:
\[
  \frac{ds}{dI} = \frac{-\mathcal{R}_0 s}{\mathcal{R}_0 s - 1}\cdot\frac{1}{N}.
\]
Integrating and using $I(\infty)=0$ yields the final size relation.
\end{proof}

\begin{codeblock}[title=\textbf{SIR model simulation}]
\begin{minted}{python}
import numpy as np
import matplotlib.pyplot as plt
from scipy.integrate import solve_ivp

N = 10000
beta, gamma = 0.3, 0.1  # R0 = 3
S0, I0, R0_init = N - 10, 10, 0

def sir(t, y):
    S, I, R = y
    dS = -beta * S * I / N
    dI = beta * S * I / N - gamma * I
    dR = gamma * I
    return [dS, dI, dR]

sol = solve_ivp(sir, [0, 200], [S0, I0, R0_init],
                t_eval=np.linspace(0, 200, 1000))

plt.figure(figsize=(10, 5))
plt.plot(sol.t, sol.y[0], 'b-', label='$S(t)$')
plt.plot(sol.t, sol.y[1], 'r-', label='$I(t)$')
plt.plot(sol.t, sol.y[2], 'g-', label='$R(t)$')
plt.xlabel('Time (days)'); plt.ylabel('Number of individuals')
plt.legend(); plt.title(f'SIR Model ($\\mathcal{{R}}_0 = {beta/gamma:.1f}$)')
plt.tight_layout(); plt.savefig('sir.pdf')
\end{minted}
\end{codeblock}

% ------------------------------------------------------------
\section{Vaccination and Herd Immunity}
\label{sec:vaccination}
% ------------------------------------------------------------

\begin{theorem}[Herd immunity threshold]
If a fraction $p$ of the population is vaccinated (and immunised),
the fraction of susceptibles is $S(0)/N = 1-p$.  The epidemic does
not start if $(1-p)\mathcal{R}_0 < 1$, i.e.:
\[
  p > p_c = 1 - \frac{1}{\mathcal{R}_0}.
\]
\end{theorem}

\begin{example}
For measles ($\mathcal{R}_0\approx 15$):
$p_c = 1 - 1/15 \approx 93.3\%$.
Over 93\% of the population must be vaccinated.
\end{example}

\begin{codeblock}[title=\textbf{Effect of vaccination}]
\begin{minted}{python}
import numpy as np
import matplotlib.pyplot as plt
from scipy.integrate import solve_ivp

N = 10000
beta, gamma = 0.3, 0.1

def sir_vacc(t, y):
    S, I, R = y
    return [-beta*S*I/N, beta*S*I/N - gamma*I, gamma*I]

fig, axes = plt.subplots(2, 2, figsize=(12, 8))
vaccination_rates = [0, 0.3, 0.6, 0.75]

for ax, p in zip(axes.flat, vaccination_rates):
    S0 = N * (1 - p) - 10
    R0_init = N * p
    I0 = 10
    sol = solve_ivp(sir_vacc, [0, 200], [S0, I0, R0_init],
                    t_eval=np.linspace(0, 200, 500))
    ax.plot(sol.t, sol.y[1], 'r-', lw=2)
    ax.set_title(f'$p = {p:.0%}$, $I_{{\\max}} = {max(sol.y[1]):.0f}$')
    ax.set_xlabel('Time'); ax.set_ylabel('$I(t)$')

plt.suptitle(f'Effect of Vaccination ($\\mathcal{{R}}_0 = {beta/gamma:.1f}$)')
plt.tight_layout(); plt.savefig('vaccination.pdf')
\end{minted}
\end{codeblock}

% ------------------------------------------------------------
\section{The SEIR Model}
\label{sec:seir}
% ------------------------------------------------------------

\begin{definition}[SEIR model]
A compartment \textbf{Exposed} ($E$) is added for individuals
infected but not yet contagious (latency period $1/\sigma$):
\begin{align}
  \frac{dS}{dt} &= -\beta\frac{SI}{N}, \\
  \frac{dE}{dt} &= \beta\frac{SI}{N} - \sigma E, \\
  \frac{dI}{dt} &= \sigma E - \gamma I, \\
  \frac{dR}{dt} &= \gamma I.
\end{align}
The basic reproduction number remains $\mathcal{R}_0 = \beta/\gamma$.
\end{definition}

\begin{codeblock}[title=\textbf{SEIR model simulation}]
\begin{minted}{python}
import numpy as np
import matplotlib.pyplot as plt
from scipy.integrate import solve_ivp

N = 10000
beta, sigma, gamma = 0.5, 0.2, 0.1
S0, E0, I0, R0_init = N - 10, 0, 10, 0

def seir(t, y):
    S, E, I, R = y
    dS = -beta * S * I / N
    dE = beta * S * I / N - sigma * E
    dI = sigma * E - gamma * I
    dR = gamma * I
    return [dS, dE, dI, dR]

sol = solve_ivp(seir, [0, 200], [S0, E0, I0, R0_init],
                t_eval=np.linspace(0, 200, 1000))

plt.figure(figsize=(10, 5))
plt.plot(sol.t, sol.y[0], 'b-', label='$S$')
plt.plot(sol.t, sol.y[1], color='orange', label='$E$')
plt.plot(sol.t, sol.y[2], 'r-', label='$I$')
plt.plot(sol.t, sol.y[3], 'g-', label='$R$')
plt.xlabel('Time'); plt.ylabel('Population')
plt.legend(); plt.title(f'SEIR Model ($\\mathcal{{R}}_0 = {beta/gamma}$)')
plt.tight_layout(); plt.savefig('seir.pdf')
\end{minted}
\end{codeblock}

% ------------------------------------------------------------
\section{Extensions and Advanced Models}
% ------------------------------------------------------------

\subsection{SIR with Demography}

\begin{definition}[SIR with demography]
Adding births and natural deaths at rate $\mu$:
\begin{align}
  \frac{dS}{dt} &= \mu N - \beta\frac{SI}{N} - \mu S, \\
  \frac{dI}{dt} &= \beta\frac{SI}{N} - (\gamma+\mu) I, \\
  \frac{dR}{dt} &= \gamma I - \mu R.
\end{align}
The basic reproduction number becomes
$\mathcal{R}_0 = \beta/(\gamma+\mu)$.
\end{definition}

\begin{proposition}[Endemic equilibrium]
For $\mathcal{R}_0 > 1$, there exists an endemic equilibrium
(disease permanently present):
\[
  S^* = \frac{N}{\mathcal{R}_0}, \qquad
  I^* = \frac{\mu N}{\gamma+\mu}\left(1-\frac{1}{\mathcal{R}_0}\right).
\]
\end{proposition}

\subsection{SIRS Model (Waning Immunity)}

If immunity is temporary (mean duration $1/\delta$), recovered
individuals become susceptible again:
\[
  \frac{dR}{dt} = \gamma I - \delta R, \qquad
  \frac{dS}{dt} = -\beta\frac{SI}{N} + \delta R.
\]
This model can produce oscillations.

% ------------------------------------------------------------
\section{Fitting to Real Data}
% ------------------------------------------------------------

\begin{codeblock}[title=\textbf{SIR fitting to data}]
\begin{minted}{python}
import numpy as np
from scipy.integrate import solve_ivp
from scipy.optimize import minimize

np.random.seed(42)
N = 10000; beta_true, gamma_true = 0.3, 0.1
sol_true = solve_ivp(lambda t, y: [-beta_true*y[0]*y[1]/N,
                     beta_true*y[0]*y[1]/N - gamma_true*y[1],
                     gamma_true*y[1]],
                     [0, 100], [N-10, 10, 0],
                     t_eval=np.arange(0, 100, 7))
I_obs = sol_true.y[1] + np.random.normal(0, 50, len(sol_true.t))
I_obs = np.maximum(I_obs, 0)

def cost(params):
    b, g = params
    if b <= 0 or g <= 0: return 1e10
    sol = solve_ivp(lambda t, y: [-b*y[0]*y[1]/N,
                    b*y[0]*y[1]/N - g*y[1], g*y[1]],
                    [0, 100], [N-10, 10, 0],
                    t_eval=np.arange(0, 100, 7))
    return np.sum((sol.y[1] - I_obs)**2)

result = minimize(cost, [0.2, 0.05], method='Nelder-Mead')
print(f"beta = {result.x[0]:.4f}, gamma = {result.x[1]:.4f}")
print(f"Estimated R0 = {result.x[0]/result.x[1]:.2f}")
\end{minted}
\end{codeblock}

% ------------------------------------------------------------
\section{Limitations of Compartmental Models}
% ------------------------------------------------------------

\begin{warning}[title=\textbf{Key limitations}]
\begin{itemize}
  \item \textbf{Homogeneity:} the population is assumed well-mixed;
    in reality, contacts are structured (social networks).
  \item \textbf{Constant parameters:} $\beta$ may vary with seasons,
    public health measures, and behaviour.
  \item \textbf{Determinism:} for small case numbers, stochastic
    models are more appropriate.
  \item \textbf{Heterogeneity:} age, geography, and comorbidities
    affect transmission and severity.
\end{itemize}
\end{warning}

% ------------------------------------------------------------
\section{Exercises}
% ------------------------------------------------------------

\begin{exercise}
For the SIR model with $N=10000$, $\beta=0.4$, $\gamma=0.1$:
\begin{enumerate}
  \item Compute $\mathcal{R}_0$ and the herd immunity threshold.
  \item Simulate the epidemic and determine the epidemic peak.
  \item Verify the final size equation numerically.
\end{enumerate}
\end{exercise}

\begin{exercise}
Compare SIR and SEIR dynamics for the same $\beta$ and $\gamma$
with $\sigma=0.2$.  How does the latency period affect the peak
and duration of the epidemic?
\end{exercise}

\begin{exercise}
Implement a stochastic SIR model (Gillespie algorithm) and compare
with the deterministic model for $N=100$ and $N=10000$.
\end{exercise}

\begin{exercise}
Model a progressive vaccination campaign: $p(t)$ increases linearly
from 0 to $p_{\max}$ over duration $T_v$.  Compare the total number
of cases for different values of $T_v$.
\end{exercise}

\begin{exercise}
Study the SIRS model with $\beta=0.5$, $\gamma=0.1$, $\delta=0.01$.
\begin{enumerate}
  \item Find the endemic equilibrium.
  \item Show numerically the existence of damped oscillations.
  \item Compare with a standard SIR model.
\end{enumerate}
\end{exercise}
